\documentclass{article}
\usepackage[utf8]{inputenc}
\usepackage{amsmath}
\usepackage{geometry}
\geometry{margin=2cm}
\begin{document}

\section*{Resolução Completa da Questão ENEM 2024 – Ciências da Natureza – Questão 128}

\subsection*{Interpretação e Estratégia}

A questão solicita identificar em qual tecido de animais vertebrados os poluentes derivados de pesticidas lipossolúveis se acumulam preferencialmente. \par

\textbf{Termos-chave:}
\begin{itemize}
  \item \textbf{Acumulação}: processo pelo qual substâncias se concentram em determinados tecidos ou órgãos.
  \item \textbf{Lipossolúvel}: substância que se dissolve em lipídios, tendendo a se acumular em tecidos gordurosos.
\end{itemize}

A estratégia envolve compreender que poluentes lipossolúveis possuem afinidade por gorduras, e portanto acumulam-se em tecidos ricos em lipídios, indicando o tecido adiposo como local preferencial. 

\subsection*{Revisão de Conteúdo e Mini Análise}

\begin{tabular}{p{0.3\textwidth} p{0.6\textwidth}}
\textbf{Conceito} & \textbf{Ideia-chave e Aplicação} \\
\hline
Lipossolubilidade & Substâncias lipossolúveis se acumulam em ambientes ricos em gordura. Poluentes apolares derivados de pesticidas acumulam-se no tecido adiposo. \\
Tecido adiposo & Rico em gordura, atua como depósito de lipídios e local preferencial de armazenamento dos poluentes lipossolúveis. \\
Distinção transporte x acúmulo & Sangue transporta substâncias mas não as armazena em concentração significativa. Poluentes circulam no sangue, porém não se acumulam nele. \\
\end{tabular}

\subsection*{Resolução e "Pulo do Gato"}

Dado que os poluentes são apolares e lipossolúveis, por suas propriedades químicas eles tendem a se dissolver e acumular em tecidos ricos em gordura. Analisando as opções, o tecido adiposo destaca-se por sua composição predominantemente lipídica, sendo o local preferencial para o depósito desses poluentes. 

\textbf{Pulo do gato:} associar diretamente o termo "lipossolúvel" ao tecido adiposo, facilitando a escolha da resposta correta.

\subsection*{Armadilhas e Distratores}

\begin{itemize}
  \item \textbf{Tecido ósseo (A):} parece um depósito estável, mas não possui lipídios suficientes para acumular poluentes lipossolúveis.
  \item \textbf{Tecido nervoso (B):} contém lipídios na bainha de mielina, mas em menor quantidade que o tecido adiposo.
  \item \textbf{Tecido epitelial (C):} contato com ambiente externo não significa alta gordura; baixo teor lipídico inviabiliza acúmulo.
  \item \textbf{Tecido sanguíneo (E):} serve para transporte, não armazenamento; poluentes não se acumulam no sangue.
\end{itemize}

\subsection*{Código da Questão}

CN_BIO_ACUMULO_001

\vspace{1em}

\textbf{Resposta correta:} 	extbf{D} – Tecido adiposo.

\begin{flushright}
\textit{Resolução endossada por análise técnica e didática.}
\end{flushright}



\end{document}